\chapter{A fresh angular budget and completion of the retained C14.1b residual}
\label{component:p4e5ba027f187}
\noindent\textbf{Component }\path{p4e5ba027f187}\par
{\small\noindent Source: \path{NEW/c1/c141b_residual_rederived.tex}\par}
{\small\noindent SHA256: \path{4e5ba027f1879415684e3c497a56aa194f0baea8216cd808d701c0b66eafa0c2}\par}
{\small\noindent Inventory use: AR1--AR21 | C14\_\allowbreak{}ENTIRE\_\allowbreak{}ORIGINAL\_\allowbreak{}PARENT\_\allowbreak{}GRAPH; precise subordinate uses retained below\par}
\medskip
\section{A fresh angular budget and completion of the retained C14.1b residual}
\label{p4e5ba027f187:sec:c141b-residual-rederived}

\paragraph{Scope and provenance.}
This is a new proof, not a restoration of the missing late-R21 source.
Keep the original physical stationary C1 source, its strict gaps,
nonnegative event coefficients, original infinitesimal total scale and
every earlier-admission complement in the surviving C14.1b statement.
Normalize the original positive coefficient to $d=1$ as in P1--P4.
Source parameters lie in a real closed extension $K_c$ and testing
angles and states in an independent extension $K_t$ of $\mathbb R$.
All conclusions are finite sums of squares in
$\mathcal T=K_c\otimes_{\mathbb R}K_t$. No order or general inverse
of a mixed tensor is used.

The exact target is the complete reciprocal-Bob polynomial P4/SP6
in the surviving \path{inherited_r19/parent_reference.tex}.
Use its notation, and set $k=\sqrt x$, $q=\sqrt z$, $z=1/L$,
$\sigma=1/\eta$, $b=Lr$, $Q=1-b$, $\theta=1-h$.
The source part of the retained residual GA9 is
\[
 x,z,\sigma>0\text{ infinitesimal},\quad
 \operatorname{res}h\ge\tfrac12,\quad k/\sigma\text{ infinitesimal},
 \quad b_0\text{ infinitesimal},\quad d_0\sim1.
 \tag{AR1}
\]
Its real testing part has $u,w>0$ infinitesimal, $WZ>0$,
$W/Z$ and $u/w$ unbounded, and $uW/(wZ)$ unbounded. These
intersections are retained. The proof below in fact needs only
$u,w,W,Z>0$ with $u,w$ infinitesimal after the inherited sign
gauges; it assumes no upper or lower bound on any state ratio.

\begin{theorem}[Fresh completion of the retained residual]
On AR1 and the exact retained testing/source intersections just
stated, the complete original parent is a finite tensor sum of
squares. Together with the surviving SF, BC, US and GA admissions
and their recorded literal, sign and projective exits, this closes
the entire original scoped C14.1b domain. It does not by itself
close C14.4 or the global C1 entry and union.
\end{theorem}

\subsection{Source coefficients and the original scalar packet}

The source estimates SF1--SF5 and US2--US3 were proved for
$z$ infinitesimal. On AR1 they imply
\[
 \begin{gathered}
 L_0:=\lambda\sigma\sim2h\text{ is a positive bounded unit},
 \qquad \lambda^{-1}\in\operatorname{Inf}_c,\\
 \rho=\lambda x/\sigma=L_0(k/\sigma)^2\in\operatorname{Inf}_c,
 \qquad \lambda x\in\operatorname{Inf}_c,
 \qquad d_G:=4D/\sigma^2\sim2. 
 \end{gathered} \tag{AR2}
\]
Here $h$ is a positive unit even when $\theta=1-h$ is
infinitesimal. In particular $\lambda>0$ and $d_G$ is a strict
original Gamma pivot. The exact sphere identity P2 proves these
bounds without an inverse of $b_0$ or $d_0$.

For clarity, the other source bounds used below are
\[
 \begin{gathered}
 P,J,M,B,\mathcal H\sim1,\quad p,l,\kappa\sim\tfrac12,
 \quad C=O(z),\quad D=O(\sigma^2),\\
 e=xh/\sigma\in\operatorname{Inf}_c,\quad
 m=\kappa xze\in\operatorname{Inf}_c,\quad
 \beta\theta\in\mathcal O_c,\quad
 \nu,K_A,\chi,R_x,A_y\in\mathcal O_c,
 \quad B_x=8\kappa e\in\operatorname{Inf}_c.
 \end{gathered} \tag{AR3}
\]
These are SF1--SF5, the P1 formulas and the SP9--SP10 polynomial
cancellations. In particular $K_A=4\kappa h/\sigma-\lambda$
is bounded although its residue $2-4\operatorname{res}h$ may be
zero or negative. Also
\[
 R_\gamma=-4\{2(C+D)(r-1)+2p+r-1+2pu^2\}
 \tag{AR4}
\]
is a bounded separated coefficient. None of AR2--AR4 contains
$u/w$, a state ratio, or a comparison of a source and a testing
rate. The only possible source poles in $a=Az=p/b$ and
$\ell=\alpha z=l/Q$ remain in positive packets below. The
possibly unbounded $\beta$ itself also remains in its original
nonnegative packet; no upper bound on $\beta$ is asserted.

Keep the exact roots
\[
 N_x=X+uhZ,\quad N_y=(Y-bwZ)/q,\quad
 \Gamma=\sigma Z-uwW,\quad R=wZ/q-uqW,
 \quad A_\theta=q(u^2-x)W.
 \tag{AR5}
\]
The whole angular packet of SP6 is
\[
 K_AA_\theta^2+\chi uA_\theta R+\mathscr S,
 \qquad
 \mathscr S=\lambda\sigma\{uZ-(x/\sigma)wW\}^2
       +\rho(u^2-x)(w^2-\sigma z)W^2.
 \tag{AR6}
\]
This is the exact P8 transformation, not a restriction of the state.
Define the four positive separated monomials
\[
 B_1=\sqrt{\lambda\sigma}\,uZ,\quad
 B_2=\sqrt\rho\,uwW,\quad
 B_3=\sqrt\lambda\,xqW,\quad
 B_4=\sqrt{\lambda x}\,wZ/q,
 \qquad G_4=(B_1B_2B_3B_4)^{1/4}.
 \tag{AR7}
\]
Every fractional power means the product of the two pure-field
powers. With $d_i=B_i-G_4$, the complete circuit is
\[
 \mathscr C:=\mathscr S+\lambda xR^2
       =\sum_{i=1}^4B_i^2-4G_4^2
       =\sum_{i=1}^4d_i^2+\mathcal E_4,
 \tag{AR8}
\]
where the explicit surviving US10 remainder is
\[
 \mathcal E_4=2G_4\left\{
 (\sqrt{B_1}-\sqrt{B_2})^2+
 (\sqrt{B_3}-\sqrt{B_4})^2+
 2\big((B_1B_2)^{1/4}-(B_3B_4)^{1/4}\big)^2\right\}.
 \tag{AR9}
\]
Its nonnegative monomial prefactors have separated square roots.

\subsection{An eight-square budget for the angular row}

The following elementary identity is the new step. Let $U,V,Z_*$
be positive separated monomials, let $a_*>0$ be pure source, and set
\[
 \begin{gathered}
 \Psi=1+U^2V^2+U^2Z_*^2+V^2Z_*^2-4UVZ_*,\\
 T=U^2-1,\qquad Q_*=VZ_*-U,\qquad
 R_*=Z_*-UV,\qquad G_*=U(V-Z_*).
 \end{gathered} \tag{AR10}
\]
There are two exact polynomial identities
\[
 \Psi=G_*^2+(1-VZ_*)^2+2VZ_*(U-1)^2,
 \qquad
 \Psi=T^2+Q_*^2+UTQ_*-UR_*^2-(U+1)G_*R_*.
 \tag{AR11}
\]
Consequently
\[
\boxed{\begin{aligned}
 &[2+a_*(U+1)]\Psi+U^2Q_*^2
       +\left[2U+\frac{U+1}{a_*}\right]R_*^2-T^2\\
 &\quad=(T+UQ_*)^2+2Q_*^2
       +(U+1)\left(\sqrt{a_*}G_*
                          -R_*/\sqrt{a_*}\right)^2\\
 &\qquad\quad+a_*(U+1)(1-VZ_*)^2
                     +2a_*(U+1)VZ_*(U-1)^2.
\end{aligned}} \tag{AR12}
\]
The right side consists of eight legal squares: expand each
factor $U+1$ into its two summands before taking roots. Only
radicals of positive separated monomials occur. In particular
neither a square root nor an inverse of $1+U$ is used. The
identity imposes no bound on $U,V,Z_*$ or on $a_*^{-1}$.

Apply this with
\[
 U=\frac uk,\qquad V=\frac w{\sqrt\sigma q},\qquad
 Z_* =\frac{\sqrt\sigma}{kq}\frac ZW,
 \qquad a_*=\lambda k,\qquad M_*=k^2qW.
 \tag{AR13}
\]
Each quotient is taken within a pure field before tensoring.
The only state inverse is the nonzero pure-testing $W^{-1}$.
The other inverses are positive pure-source factors. These
expressions are allowed to be unbounded inside a square root row.
Directly,
\[
 \mathscr C=\lambda M_*^2\Psi,\qquad
 A_\theta=M_*T,\qquad R=kqWQ_*,\qquad
 \Gamma=k\sqrt\sigma qWR_*.
 \tag{AR14}
\]
Multiplying AR12 by the square $M_*^2$ therefore proves
\[
\boxed{\quad
 \Delta_A:=\varepsilon_C\mathscr C+u^2R^2
                 +\varepsilon_G\Gamma^2-A_\theta^2
                         \ \in\ \Sigma\mathcal T^2,
 \quad} \tag{AR15}
\]
with the exact coefficients
\[
 \varepsilon_C=\frac2\lambda+u+k,\qquad
 \varepsilon_G=2\frac k\sigma u+
                         \frac{u+k}{\lambda\sigma}.
 \tag{AR16}
\]
Both are finite sums of positive bounded separated factors with
zero real residue. This follows from AR2, bounded $k/\sigma$,
and the independent facts $u\in\operatorname{Inf}_t$ and
$k,\lambda^{-1}\in\operatorname{Inf}_c$. It uses no bound on
$q u/w$, $uW/(wZ)$, or any normalized variable in AR13.
The identity also gives a direct expression for each of its eight
roots: multiply the eight expanded roots on the right of AR12
by $M_*$. Thus AR15 is a finite certificate, not a pointwise
inequality imported from an ordered compositum.

\subsection{Paying the signed angular term exactly once}

Choose a fixed positive real $c_*$ such that
$\operatorname{res}(K_A+c_*)>0$, possible because $K_A$ is
bounded. The exact allocation is
\[
\begin{split}
 K_AA_\theta^2+\mathscr S
  ={}&(K_A+c_*)A_\theta^2
       +(1-c_*\varepsilon_C)\mathscr C
       -c_*\varepsilon_G\Gamma^2\\
    &-(\lambda x+c_*u^2)R^2+c_*\Delta_A.
\end{split} \tag{AR17}
\]
AR17 uses AR8 and AR15 once each. The coefficient of the whole
circuit has residue one; its four $d_i$ rows remain strict.
The three AR9 rows with this same coefficient are retained
separately. The elementary positive-residue lemma makes each
of these bounded coefficients a tensor SOS; multiplying its
finite square decomposition by the corresponding row preserves
the finite SOS property. There is no mixed square root of
$1-c_*\varepsilon_C$.

\subsection{The two complete state normals and all original hard terms}

Use the exact CT7--CT7a radial extraction. Its constant base is
\[
 [4a+2z(2C+1)](N_y+bR)^2+4\ell(N_y-QR)^2.
 \tag{AR18}
\]
Since $p,l\sim1/2$ and $0<b,Q<1$, both coefficients exceed
one. Keep their positive pure-source excesses and retain
$(N_y+bR)^2+(N_y-QR)^2$ in the principal frame. Its bounded
two-by-two Gram has determinant $(b+Q)^2=1$ and positive
definite residue, including $b=0$ or $Q=0$ in residue. This is
a cap of whole positive squares, not a cap of unrelated matrix
entries. Retain exactly the nonconstant positive $a,\ell,C$
packets specified in CT7, before using its remaining coefficients.

The three exact remainders $E_{yy},E_{RR},E_{yR}$ in CT7a are
bounded infinitesimal: every source factor in their braces is
bounded by AR3, $C=O(z)$, $D=O(\sigma^2)$, and they have the
printed factor $z$ together with $u^2,w^2,\sigma$ or $D$.
No estimate involving a state coordinate or $u/w$ is present.

The two entire Gamma normal terms of SP6 are
\[
 -4\nu uN_x\Gamma
          +qwR_\gamma(N_y+bR)\Gamma.
 \tag{AR19}
\]
They are bounded infinitesimal entries on the original
$(N_x,N_y,R,\Gamma)$ frame, directly by AR3--AR4 and
$u,q,w$ infinitesimal. In this source sector the Gamma diagonal
is $d_G\Gamma^2$ with $d_G\sim2$, so no sphere-coordinate
inverse, Gamma conductor or weak Gamma pivot is required.
The five other complete cross terms are exactly
\[
 \begin{split}
 &-4uqw(2m+2D+1)N_xN_y+uqwR_xN_xR\\
 &\qquad+qwB_xN_xA_\theta+uA_yN_yA_\theta
                                      +\chi uA_\theta R.
 \end{split} \tag{AR20}
\]
Every displayed coefficient is infinitesimal and bounded by
AR3. The entire $N_yR$ term has already been included in
AR18 and CT7a. Thus both unrestricted state normals and every
remaining mixed term of SP6 have been allocated.

Use exactly $2N_x^2$ in the principal frame and retain exactly
$[2w^2+4mw^2]N_x^2$ as its positive excess. After AR17 the
principal frame is
\[
          (N_x,N_y,R,\Gamma,A_\theta,d_1,d_2,d_3,d_4).
 \tag{AR21}
\]
Its real residue is strictly positive definite: the $N_x$ pivot
is two, the radial block is AR18's clipped block, the Gamma
pivot is two, the angular pivot is $\operatorname{res}(K_A+c_*)>0$,
and the four circuit pivots are one. All other entries are
bounded infinitesimals, including the AR17 Gamma/radial costs.
The separated bounded-Gram lemma therefore proves finite SOS.

Finally restore, each exactly once, $4\beta(1+w^2)N_x^2$,
$2(2D+1)(uX_e)^2$, $[2w^2+4mw^2]N_x^2$,
$4bpzu^2\Gamma^2$, all CT7 positive source and testing
excesses, the AR9 remainder with its positive-residue coefficient,
and $c_*\Delta_A$. This restores precisely P4/SP6. The roots
need not be bounded; only the Gram entries have been bounded.

\subsection{Exact join and literal boundaries}

The surviving SF theorem covers $\operatorname{res}h<1/2$ in
C14.1b. US covers its unit-$h$, unit-$k/\sigma$ source for
all states and tests; US2 bounds $k/\sigma$ throughout the
unit-$h$ source. BC and GA cover their stated bounded-ratio
testing domains. Their residual is exactly contained in AR1
with the testing predicates printed above, now covered by
AR17--AR21. The inherited bounded-state, nonzero-unit-Alice,
remote, literal-axis and sign exits remain unchanged. Their
precise dependency index is the surviving
\path{proofs/c1_unbounded/DEPENDENCIES.txt}, together with
the introductions and concluding joins of SF, US and GA.
This is a finite union on the original scoped C14.1b source;
it does not erase any earlier-admission complement.

The residual itself has $D>0$ by AR2. Literal $D=0$ remains
covered by the US sphere endpoint, which divides by neither
$D$ nor a sphere coordinate. Literal $\beta=0$ is incompatible
with $H=h/\sigma$ unbounded here: P2 gives $H=H_0/B$ on
that face, a bounded quantity. Hence no nonempty C14.1b face
has been dropped. Arbitrarily small positive $b,Q,\theta,D$
or other original openings are kept wherever compatible with
the printed source sector. For complex states apply the real
routing independently to their real and imaginary parts;
each part satisfies an earlier exit or the new residual proof.
Restore only the original positive pure-source and pure-testing
normalizations.

\paragraph{Verification boundary.}
The companion exact checker verifies the elementary identities
AR11--AR12, all four scaling identities AR14, the restored
budget AR15--AR16, its eight separated squares, and the
once-paid AR17 allocation. The complete P4/SP6 and CT7a
identities are surviving algebraic inputs. Their source bounds,
the full-normal allocation and the quantified join are the
analytical argument above; an assertion count does not replace
these arguments. Independent review is requested before this
new theorem is recorded as an admitted replacement.
